\documentclass[a4paper, 14pt]{extarticle}

% Поля
%--------------------------------------
\usepackage{geometry}
\geometry{a4paper,tmargin=2cm,bmargin=2cm,lmargin=3cm,rmargin=1cm}
%--------------------------------------


%Russian-specific packages
%--------------------------------------
\usepackage[T2A]{fontenc}
\usepackage[utf8]{inputenc}
\usepackage[english, main=russian]{babel}

%--------------------------------------

\usepackage{textcomp}

% Красная строка
%--------------------------------------
\usepackage{indentfirst}               
%--------------------------------------             


%Graphics
%--------------------------------------
\usepackage{graphicx}
\graphicspath{ {./images/} }
\usepackage{float}%"Плавающие" картинки
%--------------------------------------

% Полуторный интервал
%--------------------------------------
\linespread{1.3}                    
%--------------------------------------

% Уменьшение расстояний между плавающими объектами
%--------------------------------------
\setlength{\floatsep}{10pt plus 2pt minus 2pt}
\setlength{\textfloatsep}{15pt plus 2pt minus 4pt}
\setlength{\intextsep}{10pt plus 2pt minus 2pt}
%--------------------------------------

%Выравнивание и переносы
%--------------------------------------
% Избавляемся от переполнений
\sloppy
% Запрещаем разрыв страницы после первой строки абзаца
\clubpenalty=10000
% Запрещаем разрыв страницы после последней строки абзаца
\widowpenalty=10000
%--------------------------------------

%Списки
\usepackage{enumitem}

\usepackage{minted}
%Подписи
\usepackage{caption} 

%Гиперссылки
% \usepackage{hyperref}

\usepackage[colorlinks, urlcolor = blue, filecolor = blue, citecolor = blue, linkcolor = black]{hyperref}


\hypersetup {
unicode=true
}

%Рисунки
%--------------------------------------
\DeclareCaptionLabelSeparator*{emdash}{~--- }
\captionsetup[figure]{labelsep=emdash,font=onehalfspacing,position=bottom}
%--------------------------------------

\usepackage{tempora}

%Листинги
%--------------------------------------
%Листинги
%--------------------------------------
\usepackage{listings}
\lstset{
  basicstyle=\ttfamily\footnotesize, 
  %basicstyle=\footnotesize\AnkaCoder,        % the size of the fonts that are used for the code
  breakatwhitespace=false,         % sets if automatic breaks shoulbd only happen at whitespace
  breaklines=true,                 % sets automatic line breaking
  captionpos=t,                    % sets the caption-position to bottom
  inputencoding=utf8,
  frame=single,                    % adds a frame around the code
  keepspaces=true,                 % keeps spaces in text, useful for keeping indentation of code (possibly needs columns=flexible)
  keywordstyle=\bf,       % keyword style
  numbers=left,                    % where to put the line-numbers; possible values are (none, left, right)
  numbersep=5pt,                   % how far the line-numbers are from the code
  xleftmargin=25pt,
  xrightmargin=25pt,
  showspaces=false,                % show spaces everywhere adding particular underscores; it overrides 'showstringspaces'
  showstringspaces=false,          % underline spaces within strings only
  showtabs=false,                  % show tabs within strings adding particular underscores
  stepnumber=1,                    % the step between two line-numbers. If it's 1, each line will be numbered
  tabsize=2,                       % sets default tabsize to 8 spaces
  title=\lstname                   % show the filename of files included with \lstinputlisting; also try caption instead of title
}
%--------------------------------------

%%% Математические пакеты %%%
%--------------------------------------
\usepackage{amsthm,amsfonts,amsmath,amssymb,amscd}  % Математические дополнения от AMS
\usepackage{mathtools}                              % Добавляет окружение multlined
\usepackage[perpage]{footmisc}


% графики
\usepackage{booktabs}        % Пакет для улучшенных таблиц
\usepackage{pgfplots}        % Пакет для построения графиков
\pgfplotsset{compat=1.17}    % Установка совместимости
\usepackage{tikz}            % Пакет для создания схем
\usetikzlibrary{positioning, arrows.meta}  % Библиотеки для TikZ

% Литература 

% переименовываем  список литературы в "список используемой литературы"
\addto\captionsrussian{\def\refname{Список литературы}}



\captionsetup[figure]{name=Рисунок, labelsep=space}


%--------------------------------------

%--------------------------------------
%			НАЧАЛО ДОКУМЕНТА
%--------------------------------------

\begin{document}

%--------------------------------------
%			ТИТУЛЬНЫЙ ЛИСТ
%--------------------------------------
\begin{titlepage}
\thispagestyle{empty}
\newpage


%Шапка титульного листа
%--------------------------------------
\vspace*{-60pt}
\hspace{-65pt}
\begin{minipage}{0.3\textwidth}
\hspace*{-20pt}\centering
\includegraphics[width=\textwidth]{emblem}
\end{minipage}
\begin{minipage}{0.67\textwidth}\small \textbf{
\vspace*{-0.7ex}
\hspace*{-6pt}\centerline{Министерство науки и высшего образования Российской Федерации}
\vspace*{-0.7ex}
\centerline{Федеральное государственное автономное образовательное учреждение }
\vspace*{-0.7ex}
\centerline{высшего образования}
\vspace*{-0.7ex}
\centerline{<<Московский государственный технический университет}
\vspace*{-0.7ex}
\centerline{имени Н.Э. Баумана}
\vspace*{-0.7ex}
\centerline{(национальный исследовательский университет)>>}
\vspace*{-0.7ex}
\centerline{(МГТУ им. Н.Э. Баумана)}}
\end{minipage}
%--------------------------------------

%Полосы
%--------------------------------------
\vspace{-25pt}
\hspace{-35pt}\rule{\textwidth}{2.3pt}

\vspace*{-20.3pt}
\hspace{-35pt}\rule{\textwidth}{0.4pt}
%--------------------------------------

\vspace{1.5ex}
\hspace{-35pt} \noindent \small ФАКУЛЬТЕТ\hspace{80pt} <<Информатика и системы управления>>

\vspace*{-16pt}
\hspace{47pt}\rule{0.83\textwidth}{0.4pt}

\vspace{0.5ex}
\hspace{-35pt} \noindent \small КАФЕДРА\hspace{50pt} <<Теоретическая информатика и компьютерные технологии>>

\vspace*{-16pt}
\hspace{30pt}\rule{0.866\textwidth}{0.4pt}

\vspace{11em}

\begin{center}
\Large {\bf Лабораторная работа №2} \\ 
\large <<Оценивание функций связи между ненаблюдаемыми одновременно параметрами>> \\
\large {по курсу <<Моделирование>>} \\

\end{center}\normalsize

\vspace{8em}

\vbox{%
\hfill%
\vbox{%
\hbox{Студент: Аринин М.А.}%
\hbox{Группа: ИУ9-82Б}%
\hbox{Преподаватель: Домрачева А.Б.}%
}%
} 

\bigskip

\vfill


\begin{center}
\textsl{Москва 2026}
\end{center}
\end{titlepage}
%--------------------------------------
%		КОНЕЦ ТИТУЛЬНОГО ЛИСТА
%--------------------------------------


\renewcommand{\ttdefault}{pcr}

\setlength{\tabcolsep}{3pt}

\newpage
\setcounter{page}{2}

\section{Постановка задачи}

В работе исследуется зависимость между уровнями глюкозы в крови у пациентов с разными исходами заболевания диабетом. Имеется датасет с 768 наблюдениями, содержащий медицинские показатели пациентов.

Цель работы --- проверить гипотезу о возможности описания стохастической зависимости между уровнями глюкозы для пациентов с разными исходами заболевания с помощью двупараметрической степенной функции связи. Для оценки параметров используются методы моментов, наименьших квадратов и нелинейного метода наименьших квадратов с последующим выбором наилучшего метода. Проверка гипотезы выполняется с помощью критерия Колмогорова-Смирнова.

\textbf{Входные данные:}
\begin{itemize}
  \item Датасет с 768 наблюдениями
  \item Столбец \texttt{Glucose} --- уровень глюкозы (мг/дл)
  \item Столбец \texttt{Outcome} --- исход заболевания (0 или 1)
\end{itemize}

%--------------------------------------------------------------------
\section{Построение математической модели}

В данной работе рассматривается задача оценки стохастической зависимости между уровнями глюкозы у пациентов с разными исходами заболевания. Ключевая особенность задачи заключается в том, что величины $\xi_0$ (уровень глюкозы при Outcome = 0) и $\xi_1$ (уровень глюкозы при Outcome = 1) не наблюдаются одновременно для одного и того же пациента --- это данные из разных групп пациентов. Поэтому классический подход построения функциональной зависимости по парным наблюдениям здесь неприменим.

\subsection{Формализация задачи}

Пусть $\xi_0$ --- случайная величина, представляющая уровень глюкозы для пациентов с одним диагнозом (Outcome = 0), а $\xi_1$ --- случайная величина для пациентов с иным исходом (Outcome = 1). Предполагается, что существует стохастическая зависимость, которую можно описать функцией:
\[
\xi_1 = \varphi(\xi_0) = \alpha \cdot \xi_0^{\beta},
\]
где $\alpha > 0$ и $\beta$ --- неизвестные параметры модели.

В предельном случае, когда стохастическая зависимость может быть описана функцией, функции распределения должны совпадать:
\[
F(\varphi(\xi_0)) = F(\xi_1).
\]

Это означает, что распределение преобразованных значений $\varphi(\xi_0) = \alpha \cdot \xi_0^{\beta}$ должно совпадать с распределением $\xi_1$.

\subsection{Методы оценки параметров}

Для оценки параметров $\alpha$ и $\beta$ используются три метода, из которых выбирается наилучший по критерию Колмогорова-Смирнова.

\textbf{Метод моментов:} После логарифмирования $\ln \xi_1 = \ln \alpha + \beta \ln \xi_0$ параметры оцениваются через выборочные моменты:
\[
\hat{\beta} = \sqrt{\frac{D[\ln \xi_1]}{D[\ln \xi_0]}}, \quad
\hat{\alpha} = \exp\left(\mathbb{E}[\ln \xi_1] - \hat{\beta} \cdot \mathbb{E}[\ln \xi_0]\right),
\]
где $\mathbb{E}[\cdot]$ и $D[\cdot]$ --- выборочное среднее и дисперсия соответственно.

\textbf{Метод наименьших квадратов (OLS):} Используется квантильное сопоставление распределений. Для квантилей $Q_{\xi_0}(p)$ и $Q_{\xi_1}(p)$ строится линейная регрессия $\ln Q_{\xi_1} = \ln \alpha + \beta \ln Q_{\xi_0}$ методом наименьших квадратов.

\textbf{Нелинейный метод наименьших квадратов (NLS):} Параметры находятся минимизацией суммы квадратов ошибок между квантилями:
\[
(\hat{\alpha}, \hat{\beta}) = \arg\min_{\alpha, \beta} \sum_{i=1}^{n} (Q_{\xi_1}(p_i) - \alpha \cdot Q_{\xi_0}(p_i)^{\beta})^2,
\]
где $p_i$ --- уровни квантилей.

Выбор наилучшего метода производится по минимальной статистике Колмогорова-Смирнова на контрольной выборке.

\subsection{Схема модели}

\begin{figure}[H]
  \centering
  \begin{tikzpicture}[
    node distance=1.5cm,
    box/.style={rectangle, draw, fill=#1!20, minimum width=2.5cm, minimum height=1cm, text centered, font=\footnotesize, text width=2.3cm},
    arrow/.style={->, >=stealth, thick}
  ]
    % Узлы верхнего уровня
    \node[box=blue] (data) {Исходные\\данные};
    
    % Разделение на группы
    \node[box=green, below left=1.8cm and 2cm of data] (xi0) {$\xi_0$\\Outcome=0};
    \node[box=red, below right=1.8cm and 2cm of data] (xi1) {$\xi_1$\\Outcome=1};
    
    % Разделение на выборки
    \node[box=yellow, below=of xi0] (train0) {Рабочая\\выборка\\$\xi_0$};
    \node[box=orange, below=of train0] (test0) {Контрольная\\выборка\\$\xi_0$};
    
    \node[box=yellow, below=of xi1] (train1) {Рабочая\\выборка\\$\xi_1$};
    \node[box=orange, below=of train1] (test1) {Контрольная\\выборка\\$\xi_1$};
    
    % Оценка параметров
    \node[box=purple, below=3.5cm of data] (params) {Оценка\\параметров\\$\hat{\alpha}, \hat{\beta}$};
    
    % Применение функции
    \node[box=cyan, below=of params] (phi) {$\varphi(\xi_0)$\\$=\hat{\alpha}\xi_0^{\hat{\beta}}$};
    
    % Функции распределения
    \node[box=pink, below=of phi] (dist) {Функции\\распределения\\$F(\xi_1)$, $F(\varphi)$};
    
    % Критерий
    \node[box=gray, below=of dist] (ks) {Критерий\\Колмогорова-\\Смирнова};
    
    % Стрелки от данных
    \draw[arrow] (data) -| (xi0);
    \draw[arrow] (data) -| (xi1);
    
    % Стрелки разделения на выборки
    \draw[arrow] (xi0) -- (train0);
    \draw[arrow] (train0) -- (test0);
    \draw[arrow] (xi1) -- (train1);
    \draw[arrow] (train1) -- (test1);
    
    % Стрелки к оценке параметров
    \draw[arrow] (train0.east) -- ++(0.5cm,0) |- (params.west);
    \draw[arrow] (train1.west) -- ++(-0.5cm,0) |- (params.east);
    
    % Стрелки к функции
    \draw[arrow] (params) -- (phi);
    \draw[arrow] (test0.east) -- ++(1.5cm,0) |- (phi.west);
    
    % Стрелки к функциям распределения
    \draw[arrow] (phi) -- (dist);
    \draw[arrow] (test1.west) -- ++(-1.5cm,0) |- (dist.east);
    
    % Стрелка к критерию
    \draw[arrow] (dist) -- (ks);
  \end{tikzpicture}
  \caption{Схема построения стохастической модели зависимости}
  \label{fig:model_scheme}
\end{figure}

\subsection{Критерий Колмогорова-Смирнова}

Для проверки гипотезы о соответствии распределений $F(\xi_1)$ и $F(\varphi(\xi_0))$ используется критерий Колмогорова-Смирнова. Статистика критерия вычисляется как максимальное отклонение между эмпирическими функциями распределения:
\[
D_n = \sup_x |F(\xi_1)(x) - F(\varphi(\xi_0))(x)|.
\]

Критическое значение для уровня значимости $\alpha = 0.05$:
\[
K_{\alpha} = \sqrt{-\frac{1}{2} \ln\left(\frac{1-\alpha}{2}\right)}.
\]

Гипотеза о соответствии распределений принимается, если $D_n < K_{\alpha}$, и отвергается в противном случае. Этот критерий также используется для выбора наилучшего метода оценки параметров --- выбирается метод с минимальной статистикой $D_n$.

%--------------------------------------------------------------------
\section{Решение вычислительных задач}

После формализации математической модели необходимо решить вычислительные задачи для оценки параметров и проверки гипотезы.

\subsection{Оценка параметров модели}

Для оценки параметров $\alpha$ и $\beta$ используются три метода на рабочей выборке. Поскольку величины $\xi_0$ и $\xi_1$ имеют разные размеры и не являются парными наблюдениями, используется квантильное сопоставление распределений.

\textbf{Метод моментов:} Вычисляются выборочные средние и дисперсии логарифмированных значений:
\[
\mathbb{E}[\ln \xi_0] = \frac{1}{n_0} \sum_{i=1}^{n_0} \ln \xi_0^{(i)}, \quad
D[\ln \xi_0] = \frac{1}{n_0} \sum_{i=1}^{n_0} (\ln \xi_0^{(i)} - \mathbb{E}[\ln \xi_0])^2,
\]
аналогично для $\xi_1$. Параметры оцениваются как:
\[
\hat{\beta} = \sqrt{\frac{D[\ln \xi_1]}{D[\ln \xi_0]}}, \quad
\hat{\alpha} = \exp\left(\mathbb{E}[\ln \xi_1] - \hat{\beta} \cdot \mathbb{E}[\ln \xi_0]\right).
\]

\textbf{Метод наименьших квадратов (OLS):} Для квантилей $Q_{\xi_0}(p_i)$ и $Q_{\xi_1}(p_i)$ строится линейная регрессия $\ln Q_{\xi_1} = \ln \alpha + \beta \ln Q_{\xi_0}$ методом наименьших квадратов, минимизируя:
\[
\sum_{i=1}^{n} (\ln Q_{\xi_1}(p_i) - \ln \alpha - \beta \ln Q_{\xi_0}(p_i))^2.
\]

\textbf{Нелинейный метод наименьших квадратов (NLS):} Параметры находятся прямой минимизацией:
\[
(\hat{\alpha}, \hat{\beta}) = \arg\min_{\alpha, \beta} \sum_{i=1}^{n} (Q_{\xi_1}(p_i) - \alpha \cdot Q_{\xi_0}(p_i)^{\beta})^2.
\]

Выбор наилучшего метода производится по минимальной статистике Колмогорова-Смирнова на контрольной выборке.

\subsection{Вычисление эмпирических функций распределения}

Для контрольной выборки вычисляются эмпирические функции распределения. Для выборки $\{x^{(1)}, \ldots, x^{(n)}\}$ функция распределения определяется как:
\[
F_n(x) = \frac{1}{n} \sum_{i=1}^{n} \mathbf{1}_{\{x^{(i)} \leq x\}},
\]
где $\mathbf{1}_{\{\cdot\}}$ --- индикаторная функция. Вычисляются три функции:
\begin{itemize}
  \item $F(\xi_0^{\text{test}})$ --- для контрольной выборки $\xi_0$
  \item $F(\xi_1^{\text{test}})$ --- для контрольной выборки $\xi_1$
  \item $F(\varphi(\xi_0^{\text{test}}))$ --- для преобразованных значений $\varphi(\xi_0^{\text{test}}) = \hat{\alpha} \cdot (\xi_0^{\text{test}})^{\hat{\beta}}$
\end{itemize}

Для визуализации функции вычисляются на равномерной сетке с шагом $\Delta = 0.2$ в диапазоне от минимума до максимума выборки.

\subsection{Вычисление статистики Колмогорова-Смирнова}

Для сравнения функций распределения $F(\xi_1^{\text{test}})$ и $F(\varphi(\xi_0^{\text{test}}))$ они интерполируются на общей сетке значений. Статистика критерия вычисляется как:
\[
D_n = \max_x |F(\xi_1^{\text{test}})(x) - F(\varphi(\xi_0^{\text{test}}))(x)|.
\]

Критическое значение для уровня значимости $\alpha = 0.05$:
\[
K_{\alpha} = \sqrt{-\frac{1}{2} \ln\left(\frac{1-\alpha}{2}\right)}.
\]

Гипотеза о соответствии распределений принимается при $D_n < K_{\alpha}$.

%--------------------------------------------------------------------
\section{Практическая реализация}

\newenvironment{longlisting}{\captionsetup{type=listing}}{}
\begin{longlisting}
\caption{Основной модуль анализа зависимости между уровнями глюкозы}
\begin{minted}[frame=single,fontsize = \footnotesize, linenos, xleftmargin = 1.5em, breaklines=true]{python}
import pandas as pd
import numpy as np
import matplotlib.pyplot as plt
from sklearn.model_selection import train_test_split
from scipy.interpolate import interp1d

# Загрузка данных
df = pd.read_csv('diabetes.csv')

# Разделение данных по Outcome
xi0 = df[df['Outcome'] == 0]['Glucose'].values
xi1 = df[df['Outcome'] == 1]['Glucose'].values

# Корректировка данных для возможности логарифмирования
xi0_corrected = xi0 + 1
xi1_corrected = xi1 + 1

# Разделение на рабочие и контрольные выборки
xi0_train, xi0_test = train_test_split(
    xi0_corrected, test_size=0.2, random_state=42
)
xi1_train, xi1_test = train_test_split(
    xi1_corrected, test_size=0.2, random_state=42
)

def moment_method(sample_x, sample_y):
    """Оценка параметров методом моментов"""
    log_x = np.log(sample_x)
    log_y = np.log(sample_y)
    
    E_log_x = np.mean(log_x)
    E_log_y = np.mean(log_y)
    D_log_x = np.var(log_x, ddof=0)
    D_log_y = np.var(log_y, ddof=0)
    
    beta = np.sqrt(D_log_y / D_log_x)
    ln_alpha = E_log_y - beta * E_log_x
    alpha = np.exp(ln_alpha)
    
    return alpha, beta

def ols_method(sample_x, sample_y):
    """Оценка параметров методом наименьших квадратов"""
    n_quantiles = min(len(sample_x), len(sample_y))
    quantiles = np.linspace(0.01, 0.99, n_quantiles)
    
    x_quantiles = np.quantile(sample_x, quantiles)
    y_quantiles = np.quantile(sample_y, quantiles)
    
    log_x = np.log(x_quantiles)
    log_y = np.log(y_quantiles)
    
    n = len(log_x)
    sum_log_x = np.sum(log_x)
    sum_log_y = np.sum(log_y)
    sum_log_x_sq = np.sum(log_x ** 2)
    sum_log_xy = np.sum(log_x * log_y)
    
    denominator = n * sum_log_x_sq - sum_log_x ** 2
    if abs(denominator) < 1e-10:
        return moment_method(sample_x, sample_y)
    
    beta = (n * sum_log_xy - sum_log_x * sum_log_y) / denominator
    ln_alpha = (sum_log_y - beta * sum_log_x) / n
    alpha = np.exp(ln_alpha)
    
    return alpha, beta

def nls_method(sample_x, sample_y):
    """Оценка параметров нелинейным методом наименьших квадратов"""
    from scipy.optimize import minimize
    
    n_quantiles = min(len(sample_x), len(sample_y))
    quantiles = np.linspace(0.01, 0.99, n_quantiles)
    
    x_quantiles = np.quantile(sample_x, quantiles)
    y_quantiles = np.quantile(sample_y, quantiles)
    
    def objective(params):
        alpha, beta = params
        y_pred = alpha * x_quantiles ** beta
        return np.sum((y_quantiles - y_pred) ** 2)
    
    try:
        alpha_init, beta_init = ols_method(sample_x, sample_y)
    except:
        alpha_init, beta_init = moment_method(sample_x, sample_y)
    
    result = minimize(objective, [alpha_init, beta_init], method='BFGS')
    alpha, beta = result.x
    
    return alpha, beta

# Оценка параметров тремя методами
alpha_mm, beta_mm = moment_method(xi0_train, xi1_train)
alpha_ols, beta_ols = ols_method(xi0_train, xi1_train)
alpha_nls, beta_nls = nls_method(xi0_train, xi1_train)

def empirical_distribution_function(sample):
    """Построение эмпирической функции распределения"""
    min_value = sample.min()
    max_value = sample.max()
    values = np.arange(min_value, max_value, 0.2)
    probs = []
    
    for value in values:
        prob = np.sum(sample <= value) / len(sample)
        probs.append(prob)
    
    return values, np.array(probs)

def calculate_ks_statistic(sample_y, alpha, beta, sample_x):
    """Вычисление статистики Колмогорова-Смирнова"""
    phi_x = alpha * sample_x ** beta
    
    y_sorted = np.sort(sample_y)
    phi_sorted = np.sort(phi_x)
    
    n_y = len(y_sorted)
    n_phi = len(phi_sorted)
    
    all_values = np.unique(np.concatenate([y_sorted, phi_sorted]))
    
    F_y = np.array([np.sum(y_sorted <= val) / n_y for val in all_values])
    F_phi = np.array([np.sum(phi_sorted <= val) / n_phi for val in all_values])
    
    D = np.max(np.abs(F_y - F_phi))
    return D

def calculate_mse(sample_x, sample_y, alpha, beta):
    """Вычисление среднего квадрата ошибки"""
    n_quantiles = min(len(sample_x), len(sample_y))
    quantiles = np.linspace(0.01, 0.99, n_quantiles)
    
    x_quantiles = np.quantile(sample_x, quantiles)
    y_quantiles = np.quantile(sample_y, quantiles)
    
    y_pred = alpha * x_quantiles ** beta
    mse = np.mean((y_quantiles - y_pred) ** 2)
    return mse

# Вычисление метрик для всех методов
ks_mm = calculate_ks_statistic(xi1_test, alpha_mm, beta_mm, xi0_test)
ks_ols = calculate_ks_statistic(xi1_test, alpha_ols, beta_ols, xi0_test)
ks_nls = calculate_ks_statistic(xi1_test, alpha_nls, beta_nls, xi0_test)

mse_mm = calculate_mse(xi0_test, xi1_test, alpha_mm, beta_mm)
mse_ols = calculate_mse(xi0_test, xi1_test, alpha_ols, beta_ols)
mse_nls = calculate_mse(xi0_test, xi1_test, alpha_nls, beta_nls)

# Выбор наилучшего метода
methods = [
    ("Метод моментов", mse_mm, ks_mm, alpha_mm, beta_mm),
    ("OLS", mse_ols, ks_ols, alpha_ols, beta_ols),
    ("NLS", mse_nls, ks_nls, alpha_nls, beta_nls)
]

best_method = min(methods, key=lambda x: x[2])
alpha, beta = best_method[3], best_method[4]

# Построение функций распределения для выбранного метода
values_xi0, probs_xi0 = empirical_distribution_function(xi0_test)
values_xi1, probs_xi1 = empirical_distribution_function(xi1_test)

phi_xi0 = alpha * xi0_test ** beta
values_phi, probs_phi = empirical_distribution_function(phi_xi0)

# Вычисление статистики Колмогорова-Смирнова
D_n = best_method[2]
alpha_level = 0.05
K_alpha = np.sqrt(-0.5 * np.log((1 - alpha_level) / 2))

print(f"Выбранный метод: {best_method[0]}")
print(f"α = {alpha:.6f}, β = {beta:.6f}")
print(f"Статистика Колмогорова-Смирнова: D_n = {D_n:.4f}")
print(f"Критическое значение: K_alpha = {K_alpha:.4f}")
print(f"Гипотеза {'принимается' if D_n < K_alpha else 'отвергается'}")
\end{minted}
\end{longlisting}

%--------------------------------------------------------------------
\section{Результаты}

В результате выполнения программы были оценены параметры модели тремя методами. Наилучший метод выбран по минимальной статистике Колмогорова-Смирнова на контрольной выборке.

\begin{figure}[H]
  \centering
  \includegraphics[width=1\textwidth]{methods_comparsion.png}
  \caption{Сравнение методов оценки параметров}
  \label{fig:methods_comparison}
\end{figure}

\begin{table}[H]
\centering
\begin{tabular}{|l|c|c|c|c|}
\hline
\textbf{Метод} & $\hat{\alpha}$ & $\hat{\beta}$ & MSE & $D_n$ \\
\hline
Метод моментов & 0.303030 & 1.306184 & 229.56 & 0.1100 \\
\hline
OLS (квантили) & 1.420479 & 0.979978 & 61.41 & 0.0993 \\
\hline
NLS (квантили) & 1.795230 & 0.930637 & 55.37 & 0.1204 \\
\hline
\end{tabular}
\caption{Сравнение методов оценки параметров}
\label{tab:methods_comparison}
\end{table}

\begin{figure}[H]
  \centering
  \includegraphics[width=0.7\textwidth]{functions.png}
  \caption{Функции распределения $\xi_0$ и $\xi_1$}
  \label{fig:distributions}
\end{figure}

\vspace{-0.5cm}

\begin{figure}[H]
  \centering
  \includegraphics[width=0.7\textwidth]{compare.png}
  \caption{Сравнение $F(\xi_1)$ и $F(\varphi(\xi_0))$}
  \label{fig:ks_comparison}
\end{figure}

%--------------------------------------------------------------------
\section{Анализ модели и проверка качества}

На основе полученных результатов проведён анализ качества модели.

\textbf{Анализ параметров модели:}

Выбранный метод OLS дал следующие оценки параметров: $\hat{\alpha} = 1.420479$, $\hat{\beta} = 0.979978$. 

Интерпретация параметров:
\begin{itemize}
  \item Параметр $\hat{\alpha} = 1.420479$ определяет масштаб преобразования. Значение больше единицы указывает на то, что преобразованные значения $\varphi(\xi_0)$ в среднем выше исходных значений $\xi_0$.
  \item Параметр $\hat{\beta} = 0.979978$ близок к единице, что означает почти линейную зависимость между уровнями глюкозы для пациентов с разными исходами заболевания. При $\hat{\beta} < 1$ зависимость замедляется с ростом $\xi_0$, но в данном случае замедление незначительно.
\end{itemize}

\textbf{Анализ функций распределения:}

Сравнение функций распределения $F(\xi_1)$ и $F(\varphi(\xi_0))$ на графике показывает:
\begin{itemize}
  \item Визуально кривые близки друг к другу в диапазоне значений от 75 до 200 мг/дл
  \item Максимальное расхождение наблюдается в области средних значений (около 120--150 мг/дл)
  \item В целом распределения совпадают по форме, что указывает на адекватность модели
\end{itemize}

\textbf{Проверка гипотезы критерием Колмогорова-Смирнова:}

Статистика критерия для выбранного метода OLS: $D_n = 0.0993$.

Критическое значение при уровне значимости $\alpha = 0.05$: $K_{\alpha} = 0.6101$.

\textbf{Вывод о качестве модели:}

Поскольку $D_n = 0.0993 < K_{\alpha} = 0.6101$, гипотеза о соответствии распределений принимается на уровне значимости 0.05. Это означает, что модель адекватно описывает стохастическую зависимость между уровнями глюкозы для разных исходов заболевания.

\textbf{Ограничения модели:}

\begin{itemize}
  \item Модель предполагает степенную зависимость, что может быть упрощением реальной связи
  \item Объём выборки ограничен (768 наблюдений), что влияет на точность оценок
  \item Модель не учитывает другие факторы, влияющие на уровень глюкозы
  \item Разделение по Outcome может не отражать полную картину зависимости
\end{itemize}

%--------------------------------------------------------------------
\section{Выводы}

\begin{itemize}
  \item Выполнена формализация задачи о стохастической зависимости между уровнями глюкозы для пациентов с разными исходами заболевания. Построена математическая модель на основе двупараметрической степенной функции связи $\xi_1 = \alpha \cdot \xi_0^{\beta}$, где $\xi_0$ и $\xi_1$ --- уровни глюкозы для пациентов с разными диагнозами.
  
  \item Решены вычислительные задачи: оценка параметров тремя методами (моментов, OLS, NLS) с использованием квантильного сопоставления распределений, построение эмпирических функций распределения, вычисление статистики критерия Колмогорова-Смирнова для выбора наилучшего метода и проверки гипотезы.
  
  \item Проведён анализ качества модели. Выбран наилучший метод оценки параметров по минимальной статистике Колмогорова-Смирнова. Проверена гипотеза о соответствии распределений $F(\xi_1)$ и $F(\varphi(\xi_0))$ с использованием критерия Колмогорова-Смирнова.
  
  \end{itemize}

\end{document}